% =====================================================================
%  Chapitre 7 - Sprint 4 : maintenance curative et predictive
% =====================================================================
\chapter{Sprint 4 : maintenance curative et predictive}
\label{chap:sprint4}

\section*{Introduction}
\addcontentsline{toc}{section}{Introduction}

Une flotte ne vaut que par sa disponibilite. Ce sprint traite du maintien
en condition operationnelle des vehicules, en deux temps : d'abord la
maintenance curative, qui organise le travail de l'atelier, puis
l'anticipation des pannes, qui vise a intervenir avant la defaillance
plutot qu'apres.

% =====================================================================
\section{Planification du sprint}

\subsection{Objectifs}

\begin{enumerate}
  \item Organiser le travail de l'atelier autour d'ordres de travail
    suivis par etat.
  \item Tracer les interventions et leur consommation de pieces.
  \item Tenir le stock de pieces de rechange et alerter sur les seuils
    bas.
  \item Conserver l'historique des changements d'etat.
  \item Estimer le risque de panne par vehicule et designer l'organe
    concerne.
\end{enumerate}

\subsection{Carnet du sprint}

\begin{table}[H]
\centering
\caption{Carnet du sprint 4}
\label{tab:backlog-s4}
\small
\begin{tabularx}{\textwidth}{@{}cp{2.3cm}Xcc@{}}
\toprule
\textbf{N\textsuperscript{o}} & \textbf{Acteur} & \textbf{Besoin exprime} &
\textbf{Prio.} & \textbf{Pts} \\ \midrule
1 & Responsable & Ouvrir et instruire un ordre de travail & Must & 8 \\
2 & Technicien & Faire avancer un ordre de travail par etat & Must & 5 \\
3 & Utilisateur & Consulter l'historique des changements d'etat & Should & 3 \\
4 & Technicien & Declarer et suivre une intervention & Must & 8 \\
5 & Technicien & Consommer ou restituer une piece de rechange & Must & 5 \\
6 & Responsable & Tenir le catalogue et reapprovisionner le stock & Must & 5 \\
7 & Technicien & Documenter un ordre de travail & Should & 3 \\
8 & Responsable & Consulter le risque de panne de la flotte & Could & 8 \\
9 & Technicien & Connaitre l'organe a risque d'un vehicule & Could & 5 \\
10 & Administrateur & Purger le dossier de maintenance & Must & 3 \\
\midrule
\multicolumn{4}{@{}r}{\textbf{Total}} & \textbf{53} \\
\bottomrule
\end{tabularx}
\end{table}

% =====================================================================
\section{Elaboration technique}

\subsection{Modelisation fonctionnelle}

\figReelle{maintenance-usecase.png}
  {Cas d'utilisation du service de maintenance}
  {fig:uc-maintenance}{0.60}

Le diagramme de la figure~\ref{fig:uc-maintenance} repartit le service en
cinq domaines. La lecture des associations revele une
\essentiel{repartition des roles conforme au metier} : le responsable
ouvre et instruit les ordres de travail, le technicien les fait avancer
et realise les interventions, l'administrateur seul purge le dossier.

Trois cas figurent hors des domaines, atteints par extension :
l'historique des statuts, les pieces en stock faible et la consommation
de pieces. \essentiel{Aucun n'est un but en soi} : on ne se connecte pas
pour consulter un historique, on y accede depuis l'ordre de travail que
l'on examine.

\begin{table}[H]
\centering
\caption{Description textuelle du cas « declarer une intervention »}
\label{tab:desc-intervention}
\small
\begin{tabularx}{\textwidth}{@{}p{2.5cm}X@{}}
\toprule
\textbf{Rubrique} & \textbf{Contenu} \\ \midrule
Acteur & Technicien \\
Precondition & Un ordre de travail existe et n'est pas cloture \\
Scenario nominal & Saisie du type d'intervention, de la description et du
  kilometrage ; enregistrement au statut EN\_ATTENTE ; rattachement a
  l'ordre de travail \\
Extension & Consommation de pieces : chaque piece prelevee decremente le
  stock et alimente le cout de l'intervention \\
Postcondition & L'intervention est visible dans le dossier du vehicule et
  son cout s'agrege a celui de l'ordre de travail \\
\bottomrule
\end{tabularx}
\end{table}

\subsection{Conception statique}

\figReelle{maintenance-class-domain.png}
  {Modele du domaine du service de maintenance}
  {fig:dom-maintenance}{1.00}

Le modele distingue deux notions que l'usage confond souvent.
\essentiel{L'ordre de travail} est une decision de gestion : il ouvre un
dossier, porte une priorite et un etat. \essentiel{L'intervention} est un
acte technique : elle a une duree, un cout et consomme des pieces. Un
ordre de travail peut donner lieu a plusieurs interventions.

Deux mecanismes de tracabilite le completent. \exemple{HistoriqueStatutOT}
conserve chaque changement d'etat avec son auteur et son horodatage.
\exemple{InterventionPiece} porte la quantite et le prix unitaire
\essentiel{au moment du prelevement, et non le prix courant} : le cout
d'une intervention passee ne change pas si le tarif du fournisseur
evolue.

\subsection{Conception dynamique}

\figReelle{maintenance-seq-ordre-travail.png}
  {Sequence : ouverture et avancement d'un ordre de travail}
  {fig:seq-ot}{1.00}

\figReelle{maintenance-seq-intervention.png}
  {Sequence : declaration d'une intervention avec consommation de pieces}
  {fig:seq-intervention}{1.00}

La premiere sequence montre l'ecriture systematique d'une ligne
d'historique a chaque changement d'etat. La seconde,
figure~\ref{fig:seq-intervention}, illustre l'effet de bord de la
consommation : le stock de la piece est decremente dans la meme unite de
travail que l'enregistrement de l'intervention. \essentiel{Les deux
reussissent ou echouent ensemble}, ce qui interdit qu'une piece soit
comptee consommee sans que l'intervention existe.

\subsection{Anticipation des pannes}

Le service de maintenance predictive est le huitieme et dernier service
metier ajoute au projet. Il illustre a lui seul l'exigence d'evolutivite
formulee au chapitre~\ref{chap:analyse} : \essentiel{son ajout n'a impose
de modifier aucun service existant}.

\figReelle{predictionpanne-usecase.png}
  {Cas d'utilisation du service de maintenance predictive}
  {fig:uc-prediction}{1.00}

\figReelle{predictionpanne-sequence.png}
  {Sequence : calcul du risque de panne d'un vehicule}
  {fig:seq-prediction}{1.00}

\paragraph{Ce que fait reellement le moteur.}
Il faut ici etre precis, car le terme \exemple{prediction} prete a
confusion. \essentiel{Le moteur actuellement en service n'est pas un
modele appris : c'est un ensemble d'heuristiques de maintenance.} Il
croise quatre facteurs calcules sur les donnees reelles de la
plateforme : l'usure kilometrique par organe, le vieillissement du
vehicule, la recurrence des pannes deja constatees et le retard de
revision.

Ce choix est assume et repose sur deux arguments. D'abord, le module est
\essentiel{utile immediatement}, sans jeu d'entrainement historique, que
l'entreprise ne possede pas encore. Ensuite, et surtout, chaque
prediction est \essentiel{explicable} : le moteur designe l'organe
concerne et la raison du signalement. Or c'est precisement l'explication
qu'un exploitant attend avant d'immobiliser un vehicule.

Le moteur est declare derriere une interface, et l'enumeration de la
source d'une prediction prevoit deja la valeur correspondant a un modele
appris. \essentiel{Le remplacement se fera en une ligne d'enregistrement
de dependance}, sans toucher aux couches superieures.

% =====================================================================
\section{Realisation}

\subsection{Le pilotage de l'atelier}

Les ordres de travail sont presentes en colonnes, figure~\ref{fig:s4-ot},
une colonne par etat d'avancement : a planifier, planifie, en cours,
termine. Cette disposition n'a pas ete choisie pour l'effet visuel. Elle
repond a la maniere dont le chef d'atelier raisonne : il ne cherche pas un
ordre precis, il cherche \exemple{ce qui est en attente} et \exemple{ce
qui bloque}. Une liste triee par date ne repond pas a cette question, un
tableau par etat si. Le deplacement d'une carte d'une colonne a l'autre
declenche le changement d'etat cote serveur et \essentiel{ecrit une ligne
d'historique}, conformement a la sequence commentee plus haut.

\capture{s4-ordres-travail.jpg}
  {Ordres de travail presentes par etat d'avancement}
  {fig:s4-ot}

L'intervention, figure~\ref{fig:s4-interventions}, est le travail
reellement effectue au titre d'un ordre. Elle porte le technicien, la
duree, les pieces consommees et le cout. La distinction entre l'ordre et
l'intervention n'est pas une subtilite de modelisation : un meme ordre
peut donner lieu a plusieurs interventions successives, quand une
reparation est reprise ou completee. Confondre les deux aurait rendu
impossible le calcul du cout reel d'immobilisation d'un vehicule.

\capture{s4-interventions.jpg}
  {Interventions realisees, avec leur technicien et leur cout}
  {fig:s4-interventions}

\subsection{Le stock et l'anticipation}

Le magasin de pieces, figure~\ref{fig:s4-pieces}, affiche pour chaque
reference la quantite disponible et le \essentiel{seuil d'alerte} en
dessous duquel elle doit etre reapprovisionnee. Le stock n'est jamais
saisi a la main apres une reparation : il est decremente par
l'enregistrement de l'intervention, dans la meme unite de travail. Cette
mecanique, invisible a l'ecran, est ce qui garantit que la quantite
affichee corresponde a la realite de l'etagere.

\capture{s4-pieces-rechange.jpg}
  {Stock de pieces de rechange et seuils de reapprovisionnement}
  {fig:s4-pieces}

Le tableau de bord des risques, figure~\ref{fig:s4-prediction}, restitue
le travail du moteur decrit a la section precedente. Chaque vehicule y
porte un niveau de risque, l'organe concerne et \essentiel{la raison du
signalement} : depassement du kilometrage de revision, recurrence d'une
panne, anciennete. C'est cette derniere colonne qui fait la valeur de
l'ecran. Un score seul serait recu avec scepticisme par un exploitant qui
connait ses vehicules ; un score assorti de son motif se discute, et se
verifie.

\capture{s4-prediction-pannes.jpg}
  {Tableau de bord des risques de panne, par vehicule et par organe}
  {fig:s4-prediction}

% =====================================================================
\section{Tests de validation}

\begin{table}[H]
\centering
\caption{Tests de validation du sprint 4}
\label{tab:tests-s4}
\small
\begin{tabularx}{\textwidth}{@{}cXp{2.9cm}c@{}}
\toprule
\textbf{N\textsuperscript{o}} & \textbf{Cas teste} & \textbf{Attendu} &
\textbf{Etat} \\ \midrule
T1 & Changement d'etat d'un ordre de travail & Ligne d'historique ecrite & Reussi \\
T2 & Consommation de deux pieces & Stock decremente, cout agrege & Reussi \\
T3 & Piece descendue sous son seuil & Alerte de stock faible levee & Reussi \\
T4 & Prix fournisseur modifie apres coup & Cout de l'intervention inchange & Reussi \\
T5 & Vehicule au-dela du seuil de revision & Risque signale, organe designe & Reussi \\
T6 & Purge du dossier par l'administrateur & Elements retires, vehicule conserve & Reussi \\
\bottomrule
\end{tabularx}
\end{table}

% =====================================================================
\section*{Conclusion}
\addcontentsline{toc}{section}{Conclusion}

Ce sprint livre l'atelier complet : ordres de travail suivis par etat,
interventions tracees, stock surveille et tableau de bord des risques.
Les cinquante-trois points du carnet ont ete realises.

Deux principes le structurent. Le premier est comptable : \essentiel{le
prix d'une piece est fige au prelevement}, faute de quoi l'historique des
couts se reecrirait a chaque revision tarifaire. Le second est
methodologique : \essentiel{ce qui est presente comme une prediction doit
dire de quoi il retourne}. Un moteur de regles explicables a ete prefere
a un modele appris que les donnees disponibles n'auraient pas permis
d'entrainer serieusement, et l'interface qui les separe rend la bascule
possible sans reprise.

Le sprint suivant traite de ce que la maintenance et les missions
produisent en aval : la facturation, les alertes metier et le journal des
evenements.
